\documentclass[reqno]{amsart}
\usepackage{amsmath,amssymb,amsthm}
\usepackage{fullpage}
\usepackage[colorlinks=true,linkcolor=blue]{hyperref}

\newtheorem{theorem}{Theorem}[section]
\newtheorem{lemma}[theorem]{Lemma}
\newtheorem{proposition}[theorem]{Proposition}
\theoremstyle{remark}
\newtheorem{remark}[theorem]{Remark}
\newcommand{\Z}{\mathbb Z}
\numberwithin{equation}{section}

\begin{document}
\title{Reflected layers and a cube-root saving for distinct subset sums}
\author{}
\date{}
\begin{abstract}
Let $f(n)$ be the least possible maximum of a set of $n$ positive integers
with distinct subset sums. We give a construction proving
$f(n)\leq((9/4)^{1/3}+o(1))2^n/n^{1/3}$.
A directed graph arranged in overlapping reflected layers gives modular
labels whose equal-cardinality subset sums are distinct. Its normalized
determinant is a return probability of a finite reversible Markov chain.
An explicit Green function bounds that return probability. Binary padding
then converts the modular labels to positive integers.
\end{abstract}
\maketitle

\section{Statement and modular reduction}

A finite set of positive integers is \emph{sum-distinct} if all its subset
sums, including the empty sum, are distinct. Define
\[
 f(n)=\min\{\max A: |A|=n,\ A\text{ is sum-distinct}\}.
\]
Logarithms used in choosing binary padding are to base two.

\begin{theorem}\label{thm:main}
As $n\to\infty$,
\[
 f(n)\leq\left(\left(\frac94\right)^{1/3}+o(1)\right)
                 \frac{2^n}{n^{1/3}}.
\]
In particular, $f(n)=o(2^n/\log_2 n)$.
\end{theorem}

The following elementary conversion is the same one used for the dyadic
construction in \texttt{Erdos1.tex}. We include its proof so that the present
argument is self-contained.

\begin{lemma}[Binary padding]\label{lem:padding}
Suppose $C\subseteq\{0,\ldots,q-1\}$ has size $r$ and its equal-cardinality
subset sums are distinct modulo $q$. For any integer $t\geq0$, put
\[
 b=(2^t+r)q,\qquad
 A=\{2^iq:0\leq i<t\}\cup\{b+c:c\in C\}.
\]
Then $A$ is sum-distinct, $|A|=r+t$, and
\[
 \frac{\max A}{2^{r+t}}
 <\left(1+\frac{r+1}{2^t}\right)\frac{q}{2^r}.
\]
\end{lemma}

\begin{proof}
The shifted labels exceed every binary element and are pairwise distinct.
A collision would give, with $\alpha_i,\beta_c\in\{-1,0,1\}$,
\[
 q\sum_{i<t}\alpha_i2^i+b\sum_{c\in C}\beta_c
       +\sum_{c\in C}\beta_c c=0.
\]
The absolute value of the first and third terms together is at most
$(2^t-1)q+r(q-1)<b$. Thus $\sum_c\beta_c=0$. Reduction modulo $q$ and
the hypothesis on $C$ give $\beta=0$; uniqueness of binary expansion then
gives $\alpha=0$. Finally, $\max A<(2^t+r+1)q$.
\end{proof}

\section{The graph and its lattice}

Fix integers $K\geq2$ and $s\geq1$. Write
\[
 w=K-1,\qquad L=2s,\qquad r=Lw+1.
\]
The vertices are $0,1,\ldots,Lw$, with $a=0$ and $d=Lw$.
For $0\leq h<L$ and $0\leq j\leq w$, use the notation
\[
 v_{h,j}=hw+j.
\]
Thus adjacent rows overlap in one vertex:
$v_{h,w}=v_{h+1,0}$. For $0\leq j<w$, give $v_{h,j}$ the two outgoing edges
\begin{equation}\label{eq:edges}
 v_{h,j}\longrightarrow v_{h,j+1},\qquad
 v_{h,j}\longrightarrow
 \begin{cases}
 v_{h+1,w-j},&h<L-1,\\
 a,&h=L-1.
 \end{cases}
\end{equation}
Give $d$ the single edge $d\to a$.
Every vertex except $d$ occurs exactly once as a source in
\eqref{eq:edges}; the overlapping row endpoints are not additional vertices.

The edges $i\to i+1$, together with $d\to a$, form a Hamilton cycle.
Every other edge either returns to $a$ or increases the vertex index by
$2(w-j)\geq2$. In particular, deletion of $a$ makes the graph acyclic.
Let $E$ be its adjacency matrix, with outgoing edges in rows, and put
\[
 M=2I_r-E,\qquad q=\det M,\qquad \mathcal L=\Z^rM.
\]
All row sums of $M$ are zero except the row at $d$, whose sum is one.

\begin{lemma}[Propagation through the rows]\label{lem:propagation}
There is no set $S$ of vertices satisfying all the following conditions:
\begin{enumerate}
\item $S$ contains a directed cycle and does not contain $d$;
\item every vertex in $S$ has a predecessor in $S$;
\item two distinct predecessors in $S$ force their successor to lie in $S$.
\end{enumerate}
\end{lemma}

\begin{proof}
Every cycle contains $a$, and a cycle avoiding $d$ returns to $a$ from the
last row. Between its departure and return it follows strictly increasing
indices. Such a path meets every closed row interval $[hw,(h+1)w]$:
an edge starting below $hw$ cannot end above $(h+1)w$.
Consequently $S$ meets every row.

In row zero, every vertex other than $a$ has just its immediately preceding
vertex as a predecessor. Condition (2) therefore makes the positions in $S$
in this row an initial interval $\{0,\ldots,u\}$.

For a later row $h\geq1$, position $j\in\{1,\ldots,w\}$ has precisely the
two distinct predecessors
\begin{equation}\label{eq:predecessors}
 v_{h,j-1}\quad\text{and}\quad v_{h-1,w-j}.
\end{equation}
Position zero is the last vertex of the preceding row.
Suppose first that $S$ occupies an initial interval $\{0,\ldots,u\}$ in
row $h-1$. If $u=w$, condition (3) fills row $h$. Otherwise its position
zero is absent. Its first position $j$ in $S$ must, by condition (2), have
its reflected predecessor in $S$, so $w-j\leq u$. All subsequent positions
also have reflected predecessors in $S$. Applying condition (3) successively
fills positions $j,j+1,\ldots,w$. Thus row $h$ has a nonempty final interval.

Conversely, suppose row $h-1$ has a final interval $\{\ell,\ldots,w\}$.
Position zero in row $h$ is in $S$. Its reflected inputs are in $S$ exactly
at positions $1,\ldots,w-\ell$, which condition (3) fills. Beyond that
initial interval, condition (2) says that a position in $S$ must have its
immediately preceding position in $S$. Hence the positions in $S$ form an
initial interval, with no later restart after a gap.

The rows therefore alternate between initial and final intervals. Since
$L$ is even, the last row has a final interval and contains $d$, a contradiction.
\end{proof}

\begin{lemma}[Balanced ternary avoidance]\label{lem:avoidance}
If $p\in\mathcal L$, all its coordinates belong to $\{-1,0,1\}$, and
$\sum_jp_j=0$, then $p=0$.
\end{lemma}

\begin{proof}
Write $p=zM$ with $z\in\Z^r$. The row sums give $z_d=\sum_jp_j=0$.
If $z_j>0$ but no predecessor has positive coefficient, then
\[
 p_j=2z_j-\sum_{i\to j}z_i\geq2,
\]
which is impossible. The analogous assertion holds for negative coefficients.
Following predecessors within either nonempty sign class produces a cycle,
and every cycle contains $a$. Both sign classes therefore cannot occur.

If $z\ne0$, reverse signs if necessary so that $z\geq0$ and let
$S=\{j:z_j>0\}$. This set contains a cycle, avoids $d$, and satisfies
condition (2) of Lemma~\ref{lem:propagation}. Two predecessors in $S$ give
\[
 2z_j=p_j+\sum_{i\to j}z_i\geq-1+2=1,
\]
so condition (3) also holds. Lemma~\ref{lem:propagation} gives a contradiction.
Thus $z=0$ and $p=0$.
\end{proof}

\begin{lemma}[Determinant and labels]\label{lem:labels}
The integer $q$ is positive and odd. There are $r$ distinct residues modulo
$q$ whose equal-cardinality subset sums are distinct. Moreover, $q/2^r$ is
the probability that a walk starting at $a$, choosing each outgoing edge
with probability $1/2$, dies before its first return to $a$, where the
missing probability at $d$ means death.
\end{lemma}

\begin{proof}
Put $B=E/2$. The submatrix on vertices other than $a$ is strictly upper
triangular. Its geometric series gives the total return probability $R$,
and the Schur complement gives
\[
 \det(I_r-B)=1-R=\frac{q}{2^r}.
\]
All excursions are finite; following the Hamilton path to $d$ and then
dying has positive probability. Thus $q>0$.

Every cycle cover of $E$ must be a single Hamilton cycle, because every
cycle contains $a$. The increasing order of the remaining vertices makes
the displayed Hamilton cycle the only one. Hence $\det E=\pm1$, and
$M\equiv E\pmod2$ shows that $q$ is odd.

Let $c=\operatorname{adj}(M)e_a$. Then $Mc=qe_a$, while the principal minor
obtained by deleting $a$ is $2^{r-1}$, so $c_a=2^{r-1}$.
The homomorphism
\[
 \phi:\Z^r\longrightarrow\Z/q\Z,\qquad p\longmapsto pc\pmod q
\]
is surjective and vanishes on $\mathcal L$. Both its kernel and
$\mathcal L$ have index $q$, the latter by the determinant formula for the
index of a full-rank integer row lattice. Thus $\ker\phi=\mathcal L$.
Lemma~\ref{lem:avoidance} now proves the desired property of the coordinate
labels $c_j\bmod q$. Applying it to $e_i-e_j$ proves their distinctness.
\end{proof}

\section{A return-probability estimate}

Number the positions in a row by $0,\ldots,K-1$. In a row other than the
last, the walk either takes a reflected edge or continues to the shared
endpoint, which is position zero of the next row. Its transition matrix
between rows is
\begin{equation}\label{eq:P}
 P_{i0}=2^{i+1-K},\qquad
 P_{ij}=\begin{cases}
 2^{i+j-K},&1\leq j\leq K-1-i,\\
 0,&j>K-1-i.
 \end{cases}
\end{equation}
In particular, $P_{K-1,0}=1$: entering the shared endpoint simply enters
the next row at position zero, without taking another edge.
In the final row, the death probability from position $i$ is $P_{i0}/2$.
There are $L=2s$ rows, so Lemma~\ref{lem:labels} gives
\begin{equation}\label{eq:detreturn}
 \frac{q}{2^r}=\frac12(P^{2s})_{00}.
\end{equation}

\begin{proposition}\label{prop:return}
For the matrix \eqref{eq:P} and every integer $s\geq0$,
\[
 \frac12(P^{2s})_{00}
 \leq\frac1{K+1}+\frac{K}{6(s+1)}.
\]
\end{proposition}

\begin{proof}
The matrix $P$ is stochastic and reversible for
\[
 \pi_0=\frac2{K+1},\qquad \pi_i=\frac1{K+1}\quad(1\leq i<K).
\]
Detailed balance follows immediately from \eqref{eq:P}.
Let $Q=P^2$ and $A=Q^{-1}-I_K$. Direct inversion gives
\begin{align*}
 (Af)_0&=f_0-f_1,\\
 (Af)_i&=2(2f_i-f_{i-1}-f_{i+1}) &&(1\leq i<K-1),\\
 (Af)_{K-1}&=2(f_{K-1}-f_{K-2}).
\end{align*}
For completeness, write $P=GJ$, where $J$ reverses the coordinate order
and $G$ is upper triangular. Then $G^{-1}$ has diagonal
$(2,\ldots,2,1)$ and superdiagonal $(-1,\ldots,-1)$.
Multiplying $Q^{-1}=JG^{-1}JG^{-1}$ gives the displayed formulas.

The operator $A$ is self-adjoint in $\ell^2(\pi)$, and
\[
 \langle f,Af\rangle_\pi
 =\frac2{K+1}\sum_{i=0}^{K-2}(f_{i+1}-f_i)^2.
\]
Its kernel consists of constants. Thus $Q=(I_K+A)^{-1}$ has eigenvalue
$1$ on constants and all its other eigenvalues lie in $(0,1)$.
By its orthogonal spectral decomposition, the sequence
\[
 b_j=(Q^j)_{00}-\pi_0\qquad(j\geq0)
\]
is nonnegative, nonincreasing, and summable.

We compute its sum explicitly. Let $g_i=\boldsymbol{1}_{i=0}-\pi_0$, and
let $f$ be the mean-zero solution of $Af=g$. The coordinate formulas show
that
\[
 f_i-f_{i+1}=\frac{K-1-i}{K+1}\quad(0\leq i<K-1).
\]
Using $\sum_i\pi_i f_i=0$ gives
\[
 f_0=\frac{K(K-1)(2K-1)}{6(K+1)^2}.
\]
Since $(I_K-Q)(f+g)=g$, the mean-zero geometric series for $Q$ gives
\begin{align*}
 \sum_{j=0}^{\infty} b_j
 &=f_0+g_0\\
 &=\frac{K(K-1)(2K-1)}{6(K+1)^2}+\frac{K-1}{K+1}\\
 &=\frac K3-\frac{K^2+K+6}{6(K+1)^2}
 <\frac K3.
\end{align*}
Monotonicity therefore gives
$(s+1)b_s\leq\sum_{j\geq0}b_j<K/3$.
Combining this with $\pi_0=2/(K+1)$ proves the proposition.
\end{proof}

\section{Choosing the parameters}

For each integer $K\geq2$, choose
\[
 s_K=\left\lceil\frac{K^2}{3}\right\rceil,\qquad
 r_K=2s_K(K-1)+1.
\]
Lemma~\ref{lem:labels} and Proposition~\ref{prop:return} give a modular
set of size $r_K$ with
\begin{equation}\label{eq:modbound}
 \frac{q_K}{2^{r_K}}
 \leq\frac1{K+1}+\frac{K}{6(s_K+1)}
 <\frac{3}{2K}.
\end{equation}
Apply Lemma~\ref{lem:padding} with
$t_K=2\lceil\log_2 r_K\rceil$, and put $n_K=r_K+t_K$.
The resulting positive integer set has maximum $N_K$ satisfying
\[
 \frac{N_K}{2^{n_K}}
 <\left(1+O(r_K^{-1})\right)\frac{3}{2K}.
\]
The integers $n_K$ increase with $K$, and
\[
 n_K=\frac23K^3+O(K^2),\qquad
 \frac{n_{K+1}}{n_K}\longrightarrow1.
\]

For an arbitrary sufficiently large $n$, choose the largest $K$ with
$n_K\leq n$. A sum-distinct set $A$ can be enlarged by replacing it with
$\{2a:a\in A\}\cup\{1\}$. Parity proves that the new set is sum-distinct;
its maximum doubles, so its maximum divided by $2^{|A|}$ is unchanged.
Apply this enlargement $n-n_K$ times. Since
$K=(3n/2)^{1/3}(1+o(1))$, we obtain
\[
 \frac{f(n)}{2^n}
 \leq\frac{3}{2K}(1+o(1))
 =\left(\left(\frac94\right)^{1/3}+o(1)\right)n^{-1/3}.
\]
This proves Theorem~\ref{thm:main}. The argument establishes a cube-root
saving; it does not establish an upper bound of order $2^n/\sqrt n$.

\end{document}
